\chapter{Một câu mang về}
\begin{tinhhuongbox}{Tình huống | Situation}
\songngu{Nếu cuối tiết chỉ là hết giờ, bài học sẽ ở lại rất ít.}{If the end of class is just the end of time, very little of the lesson will stay.}
\songngu{Cần một câu chốt để học sinh mang ra khỏi lớp.}{You need one closing line for students to carry out of the room.}
\end{tinhhuongbox}
\begin{noitrenlop}
\songngu{Trước khi chuông cứu tất cả chúng ta, mỗi em giữ lại cho mình một câu của hôm nay.}{Before the bell rescues all of us, keep one sentence from today for yourself.}
\end{noitrenlop}
\begin{vihieuqua}
\songngu{Câu này vừa vui vừa buộc học sinh gom lại phần đáng nhớ nhất.}{This line is funny and also forces students to gather the most memorable part.}
\end{vihieuqua}
\begin{englishpocket}
\songngu{Cụm nên nhớ: \textit{Keep one sentence from today.}}{Phrase to keep: \textit{Keep one sentence from today.}}
\end{englishpocket}
\chapter{Cho tiết học một cái tên}
\begin{tinhhuongbox}{Tình huống | Situation}
\songngu{Học sinh thường nhớ tinh thần bài tốt hơn khi tự đặt lại tên cho nó.}{Students often remember the spirit of a lesson better when they rename it themselves.}
\songngu{Đặt tên là một cách chốt rất vui mà không rỗng.}{Renaming is a fun way to close a lesson without being empty.}
\end{tinhhuongbox}
\begin{noitrenlop}
\songngu{Nếu tự đặt lại tên cho tiết này sau khi đã học xong, lớp mình đặt là gì cho đỡ hiền?}{If you could rename this lesson now that we have finished it, what would you call it to make it less bland?}
\end{noitrenlop}
\begin{vihieuqua}
\songngu{Học sinh phải nhìn lại cả tiết rồi mới đặt tên được.}{Students have to look back over the whole lesson before they can rename it.}
\end{vihieuqua}
\begin{englishpocket}
\songngu{Cụm nên nhớ: \textit{What would you call it?}}{Phrase to keep: \textit{What would you call it?}}
\end{englishpocket}
\chapter{Một lời hẹn cho tiết sau}
\begin{tinhhuongbox}{Tình huống | Situation}
\songngu{Cuối giờ đẹp là cuối giờ mở được cánh cửa cho buổi học sau.}{A beautiful ending opens a door for the next lesson.}
\songngu{Không cần hẹn to, chỉ cần hẹn đúng.}{The promise does not need to be big; it only needs to be right.}
\end{tinhhuongbox}
\begin{noitrenlop}
\songngu{Tiết sau ai mang đến lớp một ví dụ thú vị nhất sẽ mở bài giúp tôi, nên nhớ là thú vị chứ không phải kỳ quặc vô ích.}{Next lesson, whoever brings the most interesting example will open the lesson for me. Please remember that interesting is not the same as uselessly weird.}
\end{noitrenlop}
\begin{vihieuqua}
\songngu{Học sinh rời lớp với một điểm chờ nhỏ, và đó là cách giữ hậu vị.}{Students leave the room with a small point of anticipation, and that is how you keep an aftertaste.}
\end{vihieuqua}
\begin{englishpocket}
\songngu{Cụm nên nhớ: \textit{Will open the lesson for me.}}{Phrase to keep: \textit{Will open the lesson for me.}}
\end{englishpocket}
\chapter{Ai hoặc điều gì đáng cảm ơn hôm nay}
\begin{tinhhuongbox}{Tình huống | Situation}
\songngu{Kết một ngày chỉ bằng mệt mỏi thì ngày mai vào lớp sẽ nặng thêm một chút.}{If you end the day only with tiredness, tomorrow's entrance will feel a little heavier.}
\songngu{Một câu cảm ơn nhỏ giúp giữ lại điểm sáng.}{One small thank-you line helps keep a bright spot.}
\end{tinhhuongbox}
\begin{noitrenlop}
\songngu{Trước khi đi, ai hoặc điều gì đáng được cảm ơn trong hôm nay, để mai mình đỡ vào lớp với vẻ mặt như sắp trả nợ?}{Before we go, who or what deserves a thank-you today, so that tomorrow we do not walk in looking like we are about to pay a debt?}
\end{noitrenlop}
\begin{vihieuqua}
\songngu{Câu này kết ngày học bằng một nấc cảm xúc sáng hơn.}{This line ends the school day on a brighter emotional step.}
\end{vihieuqua}
\begin{englishpocket}
\songngu{Cụm nên nhớ: \textit{Who or what deserves a thank-you today?}}{Phrase to keep: \textit{Who or what deserves a thank-you today?}}
\end{englishpocket}
\chapter{Khép chuyện ở đây thôi}
\begin{tinhhuongbox}{Tình huống | Situation}
\songngu{Có những cuối tiết lớp vừa va nhau chút ít, và điều quan trọng là khép lại gọn.}{Sometimes the class has a small clash near the end, and the important thing is to close it neatly.}
\songngu{Không nên để dư vị cuối ngày là sự lôi kéo dài sang hôm sau.}{You should not let the aftertaste of the day become something dragged into tomorrow.}
\end{tinhhuongbox}
\begin{noitrenlop}
\songngu{Mình khép chuyện ở đây thôi nhé, đừng mang nó về nhà rồi mai mang lại nguyên xi.}{Let us close this here, all right? Do not take it home and bring it back tomorrow exactly the same.}
\end{noitrenlop}
\begin{vihieuqua}
\songngu{Lời chốt này cho lớp một cánh cửa đi ra thay vì một cục nặng để cầm về.}{This closing line gives the class a doorway out instead of a heavy block to carry home.}
\end{vihieuqua}
\begin{englishpocket}
\songngu{Cụm nên nhớ: \textit{Let us close this here.}}{Phrase to keep: \textit{Let us close this here.}}
\end{englishpocket}